% open-logic.tex
% dikompilasi untuk menghasilkan teks lengkap dengan gaya open-logic-dev

\documentclass[../include/open-logic-part]{subfiles}

\begin{document}

\clearpage

\begin{editorial}
Berkas ini memuat semua isi yang disertakan dalam Proyek Logika Terbuka.
Catatan editorial seperti ini, jika ditampilkan, menunjukkan bahwa berkas
tersebut dikompilasi tanpa mempertimbangkan cara materi ini akan disajikan.
Jika Anda dapat membaca catatan ini, PDF tersebut agaknya \emph{tidak
disarankan} untuk digunakan mengajar atau belajar.

Proyek Logika Terbuka menyediakan banyak mekanisme untuk menghasilkan teks
yang lebih sesuai bagi pengajaran atau belajar mandiri. Sebagai contoh, secara
baku teks akan memperlakukan semua operator logis sebagai operator primitif
dan membahas semua kasus bagi semua operator dalam pembuktian. Namun, jauh
lebih baik jika sebagian kasus itu dijadikan latihan. Proyek Logika Terbuka
juga masih terus dikembangkan. Untuk mendorong kolaborasi dan perbaikan,
materi tetap disertakan meskipun masih berupa draf, belum memiliki latihan,
dan sebagainya. PDF yang dihasilkan untuk suatu mata kuliah akan mengecualikan
bagian-bagian tersebut.

Untuk menemukan PDF yang lebih sesuai bagi pengajaran dan pembelajaran,
lihatlah \href{http://builds.openlogicproject.org/}{contoh mata kuliah yang
  tersedia di situs web OLP}. Untuk membuat versi Anda sendiri, Anda dapat
memulai dari
\href{https://github.com/OpenLogicProject/OpenLogic/tree/master/courses/sample}{berkas
  penggerak contoh} atau menelaah sumber buku-buku ajar turunan untuk melihat
contoh yang lebih menarik dan lebih lanjut.
\end{editorial}

\olimport[sets-functions-relations]{sets-functions-relations-complete}

\olimport[propositional-logic]{propositional-logic}

\olimport[first-order-logic]{first-order-logic}

\olimport[model-theory]{model-theory}

\olimport[computability]{computability}

\olimport[turing-machines]{turing-machines}

\olimport[incompleteness]{incompleteness}

\olimport[second-order-logic]{second-order-logic}

\olimport[lambda-calculus]{lambda-calculus}

\olimport[many-valued-logic]{many-valued-logic}

\olimport[normal-modal-logic]{normal-modal-logic}

\olimport[applied-modal-logic]{applied-modal-logic}

\olimport[intuitionistic-logic]{intuitionistic-logic}

\olimport[counterfactuals]{counterfactuals}

\olimport[set-theory]{set-theory}

\olimport[methods]{methods}

\olimport[history]{history}

\olimport[reference]{reference}

\end{document}
